%% figures/fig-preaccel.tex -- Figure for Sec. 3.3.
%% Pulled into notes.tex with %% figures/fig-preaccel.tex -- Figure for Sec. 3.3.
%% Pulled into notes.tex with %% figures/fig-preaccel.tex -- Figure for Sec. 3.3.
%% Pulled into notes.tex with %% figures/fig-preaccel.tex -- Figure for Sec. 3.3.
%% Pulled into notes.tex with \input{figures/fig-preaccel}.
%% Needs pgfplots, loaded in the main preamble.
%%
%% Both curves solve tau_0 adot - a = -(J/m) delta(t). Integrating across the
%% delta gives a(0+) - a(0-) = -J/(m tau_0), which fixes each branch once one
%% end condition is chosen:
%%
%%   a(t) = 0 for t<0   ->  a = -(J/m tau_0) exp(t/tau_0) for t>0   (runaway)
%%   a(t) = 0 for t>0   ->  a = +(J/m tau_0) exp(t/tau_0) for t<0   (pre-acc.)
%%
%% The two are reflections of each other through the origin, which is the point
%% of drawing them on one axis. Plotted in units of J/(m tau_0).

\begin{figure}[htb]
\centering
\begin{tikzpicture}
\begin{axis}[
  width=11.2cm, height=6.0cm,
  xlabel={$\tau/\taunought$}, ylabel={$a\ \ [\,J/m\taunought\,]$},
  xmin=-3.1, xmax=3.1, ymin=-1.45, ymax=1.45,
  axis lines=middle,
  ytick={-1,1}, yticklabels={,},
  xtick={-3,-2,-1,1,2,3},
  every axis y label/.style={at={(ticklabel* cs:1.0)},anchor=south east},
  every axis x label/.style={at={(ticklabel* cs:1.0)},anchor=north west},
  clip=true, samples=90, font=\footnotesize,
]
  %% the impulse
  \draw[->,thick,gray] (axis cs:0,0) -- (axis cs:0,1.33);
  \node[anchor=west,font=\footnotesize,gray]
       at (axis cs:0.07,1.30) {$J\,\delta(\tau)$};

  %% branch 1: no pre-acceleration, runaway afterwards
  \addplot[domain=-3.1:0,thick] {0};
  \addplot[domain=0:0.38,thick] {-exp(x)};
  \fill (axis cs:0,-1) circle (1.3pt);

  %% branch 2: asymptotic condition, pre-acceleration beforehand
  \addplot[domain=-3.1:0,thick,dashed] {exp(x)};
  \addplot[domain=0:3.1,thick,dashed] {0};
  \fill (axis cs:0,1) circle (1.3pt);

  \node[anchor=west,font=\footnotesize,align=left] at (axis cs:0.55,-0.80)
       {$a(0^-)=0$\\runs away after};
  \node[anchor=east,font=\footnotesize,align=right] at (axis cs:-0.45,0.80)
       {$a(\infty)=0$\\acts before};
\end{axis}
\end{tikzpicture}
\caption{The choice Dirac's asymptotic condition makes. Both curves solve the
same equation with the same impulse, and differ only in which end condition is
imposed. Fixing the acceleration to vanish before the blow, as an
initial-value problem would, gives the solid curve: nothing happens until
$\tau=0$ and the particle then runs away, as in Fig.~\ref{fig:runaway}(b).
Fixing it to vanish at $\tau\to+\infty$ instead gives the dashed curve: the
runaway is gone, the particle ends in uniform motion with exactly the velocity
change $J/m$ that the impulse should produce, and it acquires that velocity
\emph{before} the impulse arrives, over a time of order $\taunought$. The two
branches are reflections of each other through the origin. There is no third
branch, and no choice that avoids both.
\label{fig:preaccel}}
\end{figure}
.
%% Needs pgfplots, loaded in the main preamble.
%%
%% Both curves solve tau_0 adot - a = -(J/m) delta(t). Integrating across the
%% delta gives a(0+) - a(0-) = -J/(m tau_0), which fixes each branch once one
%% end condition is chosen:
%%
%%   a(t) = 0 for t<0   ->  a = -(J/m tau_0) exp(t/tau_0) for t>0   (runaway)
%%   a(t) = 0 for t>0   ->  a = +(J/m tau_0) exp(t/tau_0) for t<0   (pre-acc.)
%%
%% The two are reflections of each other through the origin, which is the point
%% of drawing them on one axis. Plotted in units of J/(m tau_0).

\begin{figure}[htb]
\centering
\begin{tikzpicture}
\begin{axis}[
  width=11.2cm, height=6.0cm,
  xlabel={$\tau/\taunought$}, ylabel={$a\ \ [\,J/m\taunought\,]$},
  xmin=-3.1, xmax=3.1, ymin=-1.45, ymax=1.45,
  axis lines=middle,
  ytick={-1,1}, yticklabels={,},
  xtick={-3,-2,-1,1,2,3},
  every axis y label/.style={at={(ticklabel* cs:1.0)},anchor=south east},
  every axis x label/.style={at={(ticklabel* cs:1.0)},anchor=north west},
  clip=true, samples=90, font=\footnotesize,
]
  %% the impulse
  \draw[->,thick,gray] (axis cs:0,0) -- (axis cs:0,1.33);
  \node[anchor=west,font=\footnotesize,gray]
       at (axis cs:0.07,1.30) {$J\,\delta(\tau)$};

  %% branch 1: no pre-acceleration, runaway afterwards
  \addplot[domain=-3.1:0,thick] {0};
  \addplot[domain=0:0.38,thick] {-exp(x)};
  \fill (axis cs:0,-1) circle (1.3pt);

  %% branch 2: asymptotic condition, pre-acceleration beforehand
  \addplot[domain=-3.1:0,thick,dashed] {exp(x)};
  \addplot[domain=0:3.1,thick,dashed] {0};
  \fill (axis cs:0,1) circle (1.3pt);

  \node[anchor=west,font=\footnotesize,align=left] at (axis cs:0.55,-0.80)
       {$a(0^-)=0$\\runs away after};
  \node[anchor=east,font=\footnotesize,align=right] at (axis cs:-0.45,0.80)
       {$a(\infty)=0$\\acts before};
\end{axis}
\end{tikzpicture}
\caption{The choice Dirac's asymptotic condition makes. Both curves solve the
same equation with the same impulse, and differ only in which end condition is
imposed. Fixing the acceleration to vanish before the blow, as an
initial-value problem would, gives the solid curve: nothing happens until
$\tau=0$ and the particle then runs away, as in Fig.~\ref{fig:runaway}(b).
Fixing it to vanish at $\tau\to+\infty$ instead gives the dashed curve: the
runaway is gone, the particle ends in uniform motion with exactly the velocity
change $J/m$ that the impulse should produce, and it acquires that velocity
\emph{before} the impulse arrives, over a time of order $\taunought$. The two
branches are reflections of each other through the origin. There is no third
branch, and no choice that avoids both.
\label{fig:preaccel}}
\end{figure}
.
%% Needs pgfplots, loaded in the main preamble.
%%
%% Both curves solve tau_0 adot - a = -(J/m) delta(t). Integrating across the
%% delta gives a(0+) - a(0-) = -J/(m tau_0), which fixes each branch once one
%% end condition is chosen:
%%
%%   a(t) = 0 for t<0   ->  a = -(J/m tau_0) exp(t/tau_0) for t>0   (runaway)
%%   a(t) = 0 for t>0   ->  a = +(J/m tau_0) exp(t/tau_0) for t<0   (pre-acc.)
%%
%% The two are reflections of each other through the origin, which is the point
%% of drawing them on one axis. Plotted in units of J/(m tau_0).

\begin{figure}[htb]
\centering
\begin{tikzpicture}
\begin{axis}[
  width=11.2cm, height=6.0cm,
  xlabel={$\tau/\taunought$}, ylabel={$a\ \ [\,J/m\taunought\,]$},
  xmin=-3.1, xmax=3.1, ymin=-1.45, ymax=1.45,
  axis lines=middle,
  ytick={-1,1}, yticklabels={,},
  xtick={-3,-2,-1,1,2,3},
  every axis y label/.style={at={(ticklabel* cs:1.0)},anchor=south east},
  every axis x label/.style={at={(ticklabel* cs:1.0)},anchor=north west},
  clip=true, samples=90, font=\footnotesize,
]
  %% the impulse
  \draw[->,thick,gray] (axis cs:0,0) -- (axis cs:0,1.33);
  \node[anchor=west,font=\footnotesize,gray]
       at (axis cs:0.07,1.30) {$J\,\delta(\tau)$};

  %% branch 1: no pre-acceleration, runaway afterwards
  \addplot[domain=-3.1:0,thick] {0};
  \addplot[domain=0:0.38,thick] {-exp(x)};
  \fill (axis cs:0,-1) circle (1.3pt);

  %% branch 2: asymptotic condition, pre-acceleration beforehand
  \addplot[domain=-3.1:0,thick,dashed] {exp(x)};
  \addplot[domain=0:3.1,thick,dashed] {0};
  \fill (axis cs:0,1) circle (1.3pt);

  \node[anchor=west,font=\footnotesize,align=left] at (axis cs:0.55,-0.80)
       {$a(0^-)=0$\\runs away after};
  \node[anchor=east,font=\footnotesize,align=right] at (axis cs:-0.45,0.80)
       {$a(\infty)=0$\\acts before};
\end{axis}
\end{tikzpicture}
\caption{The choice Dirac's asymptotic condition makes. Both curves solve the
same equation with the same impulse, and differ only in which end condition is
imposed. Fixing the acceleration to vanish before the blow, as an
initial-value problem would, gives the solid curve: nothing happens until
$\tau=0$ and the particle then runs away, as in Fig.~\ref{fig:runaway}(b).
Fixing it to vanish at $\tau\to+\infty$ instead gives the dashed curve: the
runaway is gone, the particle ends in uniform motion with exactly the velocity
change $J/m$ that the impulse should produce, and it acquires that velocity
\emph{before} the impulse arrives, over a time of order $\taunought$. The two
branches are reflections of each other through the origin. There is no third
branch, and no choice that avoids both.
\label{fig:preaccel}}
\end{figure}
.
%% Needs pgfplots, loaded in the main preamble.
%%
%% Both curves solve tau_0 adot - a = -(J/m) delta(t). Integrating across the
%% delta gives a(0+) - a(0-) = -J/(m tau_0), which fixes each branch once one
%% end condition is chosen:
%%
%%   a(t) = 0 for t<0   ->  a = -(J/m tau_0) exp(t/tau_0) for t>0   (runaway)
%%   a(t) = 0 for t>0   ->  a = +(J/m tau_0) exp(t/tau_0) for t<0   (pre-acc.)
%%
%% The two are reflections of each other through the origin, which is the point
%% of drawing them on one axis. Plotted in units of J/(m tau_0).

\begin{figure}[htb]
\centering
\begin{tikzpicture}
\begin{axis}[
  width=11.2cm, height=6.0cm,
  xlabel={$\tau/\taunought$}, ylabel={$a\ \ [\,J/m\taunought\,]$},
  xmin=-3.1, xmax=3.1, ymin=-1.45, ymax=1.45,
  axis lines=middle,
  ytick={-1,1}, yticklabels={,},
  xtick={-3,-2,-1,1,2,3},
  every axis y label/.style={at={(ticklabel* cs:1.0)},anchor=south east},
  every axis x label/.style={at={(ticklabel* cs:1.0)},anchor=north west},
  clip=true, samples=90, font=\footnotesize,
]
  %% the impulse
  \draw[->,thick,gray] (axis cs:0,0) -- (axis cs:0,1.33);
  \node[anchor=west,font=\footnotesize,gray]
       at (axis cs:0.07,1.30) {$J\,\delta(\tau)$};

  %% branch 1: no pre-acceleration, runaway afterwards
  \addplot[domain=-3.1:0,thick] {0};
  \addplot[domain=0:0.38,thick] {-exp(x)};
  \fill (axis cs:0,-1) circle (1.3pt);

  %% branch 2: asymptotic condition, pre-acceleration beforehand
  \addplot[domain=-3.1:0,thick,dashed] {exp(x)};
  \addplot[domain=0:3.1,thick,dashed] {0};
  \fill (axis cs:0,1) circle (1.3pt);

  \node[anchor=west,font=\footnotesize,align=left] at (axis cs:0.55,-0.80)
       {$a(0^-)=0$\\runs away after};
  \node[anchor=east,font=\footnotesize,align=right] at (axis cs:-0.45,0.80)
       {$a(\infty)=0$\\acts before};
\end{axis}
\end{tikzpicture}
\caption{The choice Dirac's asymptotic condition makes. Both curves solve the
same equation with the same impulse, and differ only in which end condition is
imposed. Fixing the acceleration to vanish before the blow, as an
initial-value problem would, gives the solid curve: nothing happens until
$\tau=0$ and the particle then runs away, as in Fig.~\ref{fig:runaway}(b).
Fixing it to vanish at $\tau\to+\infty$ instead gives the dashed curve: the
runaway is gone, the particle ends in uniform motion with exactly the velocity
change $J/m$ that the impulse should produce, and it acquires that velocity
\emph{before} the impulse arrives, over a time of order $\taunought$. The two
branches are reflections of each other through the origin. There is no third
branch, and no choice that avoids both.
\label{fig:preaccel}}
\end{figure}
